\documentclass[11pt,a4paper]{article}

\usepackage[utf8]{inputenc}
\usepackage[T1]{fontenc}
\usepackage[italian]{babel}
\usepackage{lmodern}
\usepackage{amsmath}
\usepackage{booktabs}
\usepackage{geometry}
\usepackage{listings}
\usepackage{xcolor}
\usepackage{enumitem}
\usepackage[hidelinks]{hyperref}

\geometry{margin=2.6cm}

% stile per i blocchi di codice
\definecolor{codebg}{RGB}{245,245,245}
\definecolor{codeframe}{RGB}{210,210,210}

\lstdefinestyle{plain}{
  backgroundcolor=\color{codebg},
  frame=single,
  rulecolor=\color{codeframe},
  basicstyle=\ttfamily\footnotesize,
  numbers=left,
  numberstyle=\tiny\color{gray},
  numbersep=8pt,
  xleftmargin=14pt,
  xrightmargin=6pt,
  breaklines=true,
  showstringspaces=false,
  columns=fullflexible,
  keepspaces=true,
  upquote=true
}
\lstdefinestyle{py}{
  style=plain,
  language=Python,
  keywordstyle=\color{blue!60!black},
  commentstyle=\color{green!45!black},
  stringstyle=\color{red!55!black}
}
\lstset{style=plain}

\title{\vspace{-1.5cm}\textbf{PROGETTO MAADB --- LAB 2025/2026}\\[0.3cm]
\large FinDetect: polyglot persistence Neo4j + MongoDB\\ sul benchmark LDBC FinBench}
\author{Pietro Sardi (matr. 989803) \quad Pietro Longo (matr. 1062127)\\[0.2cm]
\normalsize Universit\`a degli Studi di Torino}
\date{A.A. 2025--2026}

\begin{document}

\maketitle
\thispagestyle{empty}

\begin{abstract}
\noindent FinDetect \`e un prototipo di applicazione web per l'interrogazione e
l'analisi del dataset finanziario \emph{LDBC FinBench}. Gli stessi dati sono
caricati su due basi di dati non relazionali --- un database a grafo (Neo4j) e
un database a documenti (MongoDB) --- per sfruttare ciascun motore dove rende di
pi\`u: il grafo per le interrogazioni sui cammini fra conti, il documento per le
schede anagrafiche aggregate. L'app espone sei query parametriche (lookup e
analitiche, alcune cross-database) tramite un'interfaccia Flask. La relazione
discute le scelte progettuali, le collega ai concetti del corso e analizza in
modo critico i tempi di esecuzione misurati sui dati reali.
\end{abstract}

\newpage
\tableofcontents
\newpage

% ===================================================================
\section{Introduzione}
% ===================================================================

Il problema che affrontiamo \`e l'analisi di dati finanziari per scopi di
\emph{antifrode}: dato un grande insieme di conti, persone, aziende e
transazioni, vogliamo rispondere a domande sia puntuali (``mostrami il profilo
di questa persona'') sia analitiche (``trova i conti che fanno da snodo a
tantissime transazioni'', ``trova i giri di denaro che tornano al mittente'').

Il dato di partenza \`e il benchmark \textbf{LDBC FinBench}, uno standard pensato
proprio per stressare i sistemi di gestione dati su carichi di lavoro di tipo
finanziario. La traccia del progetto chiede di caricare gli stessi dati su due
sistemi non relazionali e di costruire una web-app che permetta query
parametriche di tipo \emph{lookup} e query \emph{analitiche}, di cui almeno una
per categoria deve essere cross-database.

La scelta di non usare un singolo database relazionale nasce da
un'osservazione che \`e anche al centro del programma del corso: il modello
relazionale \`e ottimo quando i dati sono ``piatti e omogenei'', ma diventa
scomodo quando l'informazione interessante sta nelle \emph{relazioni fra le
entit\`a} (cfr. appunti del corso, cap. 1.2 e 1.4). Trovare un ciclo di
pagamenti $A\to B\to C\to A$ in SQL richiederebbe pi\`u self-join annidati e
costosi; su un grafo \`e una singola interrogazione di
\emph{raggiungibilit\`a}. Da qui la scelta della \emph{polyglot persistence}:
usare due tecnologie diverse, ognuna nel suo dominio di forza.

\paragraph{Cosa fa l'applicazione.} L'utente sceglie una delle sei query da
un men\`u, imposta i parametri (con la possibilit\`a di pescarne di realistici a
caso dal dataset) e ottiene una tabella di risultati con il numero di righe e il
tempo di esecuzione. Il prototipo \`e volutamente minimale, come richiesto dalla
traccia: l'obiettivo non \`e un prodotto, ma una dimostrazione funzionante delle
scelte di modellazione.

% ===================================================================
\section{Accenni teorici}
% ===================================================================

Questa sezione richiama i concetti che useremo per giustificare le scelte. Ogni
nozione \`e spiegata in modo intuitivo e, dove possibile, collegata agli appunti
del corso.

\subsection{Modello a grafo e modello a documenti}

Un \textbf{database a grafo} rappresenta i dati come nodi (le entit\`a) e archi
(le relazioni), entrambi dotati di propriet\`a. A differenza del vecchio modello
reticolare, qui il grafo non \`e imposto ovunque ma \emph{scelto dove la natura
dei dati \`e gi\`a relazionale} (cfr. appunti, cap. 1.4, ``Basi di dati a
grafo''). La forza del modello \`e che abilita le \textbf{query di
raggiungibilit\`a}, cio\`e interrogazioni che chiedono se e come due nodi sono
collegati da un cammino, che nel relazionale ``restano impossibili'' o molto
costose. \`E esattamente il caso dei nostri cicli di transazioni (Q4) e della
ricerca degli snodi (Q6).

Un \textbf{database a documenti} memorizza invece oggetti autodescrittivi (in
MongoDB, documenti BSON simili a JSON). Si rinuncia alla rigidit\`a dello schema
relazionale in cambio di flessibilit\`a: un documento pu\`o contenere
direttamente le sue sotto-strutture annidate. Questo \`e ideale per la
\emph{scheda di una persona}, dove vogliamo anagrafica, conti e prestiti gi\`a
raccolti insieme senza dover ricomporre tutto con dei join (Q2).

\paragraph{Perch\'e due database e non uno.} La \emph{polyglot persistence}
\`e la pratica di usare pi\`u tecnologie di memorizzazione nello stesso sistema,
ciascuna per il tipo di accesso a cui \`e pi\`u adatta. Il costo \`e la
\textbf{ridondanza controllata}: alcune informazioni (es. i transfer) vivono sia
come archi in Neo4j sia come documenti in MongoDB. Lo accettiamo
consapevolmente perch\'e i due motori servono interrogazioni diverse: il grafo
per i cammini, il documento per le aggregazioni anagrafiche. \`E un classico
esempio di scelta di ingegneria che baratta un po' di spazio e di coerenza in
cambio di prestazioni mirate.

\paragraph{Perch\'e proprio Neo4j e MongoDB.} Fra i database a grafo abbiamo
scelto \textbf{Neo4j} perch\'e \`e il pi\`u maturo e diffuso, ha un'edizione
Community gratuita e usa \emph{Cypher}, un linguaggio dichiarativo molto
leggibile (un pattern come $A\to B\to C\to A$ si scrive quasi come si
disegnerebbe su carta). Alternative come ArangoDB o Amazon Neptune sarebbero
state valide ma meno immediate da installare e usare in locale. Fra i document
store abbiamo scelto \textbf{MongoDB} per la stessa ragione di diffusione, per
la potenza della \emph{aggregation pipeline} e per la semplicit\`a con cui si
configura un replica set; un'alternativa come PostgreSQL con campi JSONB sarebbe
stata di fatto un ritorno al relazionale, fuori dallo spirito ``due sistemi non
relazionali'' richiesto dalla traccia.

\subsection{Centralit\`a di grado per la \emph{hub detection}}

La query Q6 cerca gli ``snodi'' della rete, cio\`e i conti che partecipano a un
numero anomalo di transazioni. La misura naturale \`e la \textbf{centralit\`a di
grado} (\emph{degree centrality}): il grado di un nodo \`e il numero di archi che
lo toccano. Considerando le transazioni entranti e uscenti separatamente:
%
\begin{equation}
\deg(v) \;=\; \deg_{\text{in}}(v) + \deg_{\text{out}}(v)
\;=\; \bigl|\{\,e \in E : e = (u,v)\,\}\bigr| + \bigl|\{\,e \in E : e = (v,w)\,\}\bigr|
\label{eq:degree}
\end{equation}
%
dove $E$ \`e l'insieme degli archi \texttt{TRANSFER}. Un conto con
$\deg(v)$ molto alto \`e un candidato hub: nello scenario antifrode \`e un punto
che merita attenzione. Q6 calcola questo grado \emph{dentro una finestra
temporale} e offre due criteri per decidere chi \`e hub:
\begin{itemize}[nosep]
  \item \textbf{soglia assoluta} (parametro \texttt{minDegree}, il default):
  tiene i conti con $\deg(v) \ge$ soglia. \`E il criterio \emph{operativo},
  pi\`u vicino al caso d'uso antifrode/vigilanza dichiarato da LDBC FinBench: si
  vuole segnalare chi supera una certa attivit\`a indipendentemente dal resto
  del sistema;
  \item \textbf{top $x\%$} (parametro opzionale \texttt{topPct}): tiene la
  frazione pi\`u attiva della distribuzione. \`E un criterio \emph{relativo},
  adatto ad esplorare come sono distribuiti i gradi; per calcolarlo serve
  per\`o l'intera classifica dei conti attivi, quindi costa un po' di pi\`u.
\end{itemize}
In entrambi i casi la query riporta, oltre ai primi 20 hub, quanti conti
soddisfano il criterio e il grado minimo fra loro: cos\`i i due approcci si
confrontano direttamente (ad esempio, a SF~1 sul triennio 2020--2022 il top
$1\%$ corrisponde a circa 2.600 conti con grado $\ge 72$, mentre la soglia 500
ne isola poche decine).

\subsection{Selettivit\`a e costo delle query}

Per ragionare sui tempi di esecuzione useremo il concetto di
\textbf{selettivit\`a} di una condizione, definito negli appunti (cap. 5,
ottimizzazione delle query) come il rapporto fra la cardinalit\`a del risultato e
quella dell'input:
%
\begin{equation}
\text{sel}(P) \;=\; \frac{|\sigma_P(R)|}{|R|} \in [0,1]
\label{eq:sel}
\end{equation}
%
Un valore basso significa alta selettivit\`a (il filtro lascia passare poche
tuple). Le nostre query di lookup sono molto selettive (filtrano su un singolo
\texttt{id} indicizzato) e infatti costano pochi millisecondi; le analitiche a
grafo come Q4 e Q6 sono poco selettive (devono esplorare molti archi) e questo
si rifletter\`a nei tempi misurati (Sezione~\ref{sec:risultati}).

\subsection{Aspetti distribuiti: replica, consistenza, partizionamento}
\label{sec:distr}

La traccia chiede di collegare il progetto, dove pertinente, ai concetti di
sistema distribuito (replica, partizionamento/sharding, consistenza, batch e
stream processing). \`E corretto premettere che il programma del corso, e i
relativi appunti, trattano questi temi nel contesto del DBMS \emph{centralizzato}
(propriet\`a ACID, controllo di concorrenza, recovery): la nozione di
``consistenza'' degli appunti \`e quella transazionale ACID (cap. 6), non la
consistenza distribuita. Per gli aspetti distribuiti veri e propri ci appoggiamo
quindi alla documentazione ufficiale dei sistemi usati, collegandoli a come il
prototipo \`e effettivamente configurato.

\paragraph{Replica (MongoDB replica set).} Nel nostro
\texttt{docker-compose.yml} MongoDB \`e avviato in modalit\`a \emph{replica set}
(\texttt{-{}-replSet rs0}), con uno script di inizializzazione che chiama
\texttt{rs.initiate()}. Un replica set \`e un gruppo di nodi che mantengono la
stessa copia dei dati con architettura \emph{primary-secondary}: le scritture
vanno al primario e vengono propagate ai secondari tramite un log delle
operazioni (\emph{oplog}). Nel prototipo il set ha un solo nodo --- sufficiente
in laboratorio --- ma l'architettura \`e quella replicata. \textbf{Perch\'e
averlo, anche con un nodo?} Perch\'e MongoDB abilita le transazioni
multi-documento (con le loro garanzie ACID) \emph{solo} sopra un replica set: il
set \`e quindi un prerequisito architetturale, non un dettaglio.

\paragraph{Consistenza.} Sopra un replica set MongoDB offre una consistenza
\emph{regolabile} tramite \texttt{write concern} (quanti nodi devono confermare
una scrittura) e \texttt{read concern} (quanto recente deve essere la lettura).
Con un solo nodo la lettura \`e di fatto fortemente consistente; in un set a pi\`u
nodi si potrebbe scegliere fra letture pi\`u veloci ma potenzialmente
indietro nel tempo (dai secondari) e letture pi\`u sicure (dal primario). \`E il
tipico compromesso fra latenza e consistenza dei sistemi distribuiti, assente
nel modello puramente centralizzato del corso.

\paragraph{Partizionamento/sharding.} Non lo abbiamo usato: alla scala del
nostro dataset (\emph{scale factor} 1, descritto sotto) i dati stanno
comodamente su un singolo nodo, e introdurre lo sharding sarebbe stato
\textbf{over-engineering}. Lo sharding diventa necessario quando i dati non
entrano pi\`u su una macchina: MongoDB partiziona orizzontalmente le collezioni
in base a una \emph{shard key}, Neo4j Community (quello che usiamo) non supporta
il clustering. Lo citiamo come scelta consapevolmente scartata per la scala in
gioco.

\paragraph{Batch vs stream processing.} Il dataset \`e prodotto in
\textbf{batch}: il datagen ufficiale LDBC gira su \emph{Apache Spark}, il motore
di elaborazione distribuita per eccellenza, e genera i CSV in una sola passata.
Anche il nostro ETL \`e batch (caricamento massivo dei CSV nei due database). Non
usiamo \emph{stream processing} perch\'e i dati sono statici, generati una volta
sola; un sistema antifrode reale, che deve reagire alle transazioni mentre
accadono, userebbe invece un motore a flussi (es. Kafka + Spark Streaming) con
elaborazione a finestre temporali. Lo segnaliamo come naturale estensione.

% ===================================================================
\section{Dataset e preprocessing}
% ===================================================================

\subsection{Il modello dei dati FinBench}

Il dataset descrive un piccolo mondo finanziario fatto di poche entit\`a e di una
fitta rete di relazioni fra loro. Le entit\`a (i nodi in Neo4j, le collezioni in
MongoDB) sono:
\begin{itemize}[nosep]
  \item \textbf{Person} e \textbf{Company}: i titolari, persone fisiche e aziende;
  \item \textbf{Account}: i conti, ciascuno con un tipo (es. \emph{retirement
  account}, \emph{credit card}) e un intestatario;
  \item \textbf{Loan}: i prestiti, con importo, saldo e finalit\`a;
  \item \textbf{Medium}: gli strumenti di pagamento usati nelle operazioni.
\end{itemize}
Il cuore del dataset sono per\`o le \emph{relazioni}: una persona \emph{possiede}
conti (\texttt{personOwnAccount}) e \emph{richiede} prestiti
(\texttt{personApplyLoan}); soprattutto, i conti si scambiano denaro con i
\texttt{TRANSFER} --- il tipo di arco pi\`u importante per noi, oltre 1,3 milioni
di istanze --- e fanno operazioni di \texttt{DEPOSIT}, \texttt{WITHDRAW} e
\texttt{REPAY}. \`E proprio questa rete di transazioni fra conti a rendere
naturale il modello a grafo: le nostre query analitiche (cicli, hub) sono
domande \emph{sulla forma della rete}, non sui singoli record.

\subsection{LDBC FinBench e lo \emph{scale factor}}

Il dataset \`e generato dal datagen ufficiale LDBC FinBench (tag
\texttt{v0.1.0}) con \emph{scale factor} (SF) pari a 1. Lo scale factor \`e il
``rubinetto'' che regola la dimensione dei dati generati: SF~1 produce un dataset
abbastanza grande da essere interessante (milioni di archi) ma abbastanza piccolo
da girare su un portatile.

\paragraph{Perch\'e SF~1 e non di pi\`u.} La generazione \`e l'operazione pi\`u
pesante dell'intero progetto, perch\'e Spark deve simulare l'attivit\`a
finanziaria. Su un portatile Apple Silicon la generazione di SF~1 ha richiesto
circa 7--8 minuti, limitando la memoria a 8~GB. SF pi\`u alti (10, 100) avrebbero
moltiplicato i tempi e lo spazio senza aggiungere nulla alla dimostrazione delle
scelte di modellazione, che \`e l'obiettivo del progetto. \`E la stessa logica
con cui in un esperimento si sceglie un campione rappresentativo invece
dell'intera popolazione.

\subsection{Volumi caricati (dati reali)}

La Tabella~\ref{tab:neo4j} e la Tabella~\ref{tab:mongo} riportano i conteggi
\emph{effettivamente misurati} sui due database dopo l'ETL. In Neo4j il dataset
\`e fatto di \textbf{643.241 nodi} e \textbf{5.372.701 archi}.

\begin{table}[h]
\centering
\begin{minipage}{0.48\textwidth}
\centering
\begin{tabular}{lr}
\toprule
\textbf{Nodo (label)} & \textbf{Conteggio} \\
\midrule
Account & 264.075 \\
Loan    & 159.166 \\
Medium  & 100.000 \\
Person  & 80.000 \\
Company & 40.000 \\
\midrule
\textbf{Totale nodi} & \textbf{643.241} \\
\bottomrule
\end{tabular}
\end{minipage}
\hfill
\begin{minipage}{0.48\textwidth}
\centering
\begin{tabular}{lr}
\toprule
\textbf{Arco (tipo)} & \textbf{Conteggio} \\
\midrule
WITHDRAW  & 2.011.359 \\
TRANSFER  & 1.379.527 \\
DEPOSIT   & 512.680 \\
SIGNIN    & 451.362 \\
OWN       & 264.075 \\
REPAY     & 262.571 \\
INVEST    & 260.156 \\
APPLY     & 159.166 \\
GUARANTEE & 71.805 \\
\midrule
\textbf{Totale archi} & \textbf{5.372.701} \\
\bottomrule
\end{tabular}
\end{minipage}
\caption{Volumi in Neo4j: nodi per label e archi per tipo.}
\label{tab:neo4j}
\end{table}

\begin{table}[h]
\centering
\begin{tabular}{lr lr}
\toprule
\textbf{Collezione} & \textbf{Documenti} & \textbf{Collezione} & \textbf{Documenti} \\
\midrule
account          & 264.075   & person            & 80.000 \\
transfer         & 1.379.527 & company           & 40.000 \\
loan             & 159.166   & personOwnAccount  & 175.956 \\
medium           & 100.000   & companyOwnAccount & 88.119 \\
personApplyLoan  & 106.346   &                   &        \\
\bottomrule
\end{tabular}
\caption{Volumi in MongoDB (database \texttt{findetect}): documenti per collezione.}
\label{tab:mongo}
\end{table}

Si noti la coerenza fra i due database: i 264.075 conti sono sia nodi
\texttt{Account} sia documenti \texttt{account}, e gli 1.379.527 transfer sono
sia archi \texttt{TRANSFER} sia documenti \texttt{transfer}. Anche le due
collezioni-link \texttt{personOwnAccount} e \texttt{companyOwnAccount} sommano
esattamente a 264.075: in FinBench ogni conto ha un solo intestatario, persona
\emph{o} azienda, ed \`e per questo che le query cross-database (Q3, Q6)
interrogano entrambe. \`E la ridondanza
controllata di cui sopra: la stessa entit\`a vive nei due mondi con la forma pi\`u
comoda per ciascuno.

\subsection{ETL: dal CSV ai due database}

Il caricamento \`e gestito da due script Python (\texttt{etl/load\_neo4j.py} e
\texttt{etl/load\_mongo.py}) che leggono i CSV grezzi prodotti dal datagen e li
inseriscono nei rispettivi database. Le scelte principali:

\begin{itemize}[nosep]
  \item \textbf{Caricamento a blocchi (batch).} Gli inserimenti sono raggruppati
  in lotti invece di una riga alla volta: scrivere milioni di record uno per uno
  sarebbe lentissimo per via del costo fisso di ogni operazione di I/O.
  \item \textbf{Indici prima dei dati.} In Neo4j creiamo un indice su
  \texttt{Account.id} prima di caricare gli archi, cos\`i la ricerca dei nodi a
  cui agganciare ogni \texttt{TRANSFER} non degenera in una scansione completa.
  In MongoDB l'indice su \texttt{\_id} \`e automatico.
  \item \textbf{Le date come interi.} I timestamp sono salvati come millisecondi
  dall'epoch (interi), pi\`u compatti e veloci da confrontare delle stringhe.
\end{itemize}

% ===================================================================
\section{Architettura del sistema}
% ===================================================================

\subsection{Struttura del codice}

Il progetto \`e organizzato in moduli piccoli e con un compito chiaro:

\begin{itemize}[nosep]
  \item \texttt{docker-compose.yml} --- avvia Neo4j e MongoDB in container;
  \item \texttt{datagen/run\_datagen.sh} --- clona, compila e lancia il datagen Spark;
  \item \texttt{etl/} --- gli script di caricamento CSV verso i due database;
  \item \texttt{app/db/} --- due moduli (\texttt{mongo.py}, \texttt{neo4j.py})
  che incapsulano la connessione: il resto del codice chiede una connessione e
  non sa nulla di come viene creata;
  \item \texttt{app/queries/} --- un file per ognuna delle sei query, pi\`u un
  modulo \texttt{\_common.py} con le funzioni condivise (conversione date,
  misura dei tempi);
  \item \texttt{app/app.py} --- l'applicazione Flask con le rotte.
\end{itemize}

\paragraph{Perch\'e separare connessione e query.} Ogni modulo in
\texttt{app/db/} espone una sola funzione (\texttt{get\_driver()},
\texttt{get\_db()}) che restituisce una connessione riusabile. \`E una forma
leggera di \emph{repository pattern}: se domani cambiassimo le credenziali o
spostassimo un database, toccheremmo un solo file invece di tutte le query.

\subsection{Le sei query e la loro classificazione}

La traccia chiede almeno due query di lookup (di cui una cross-database) e almeno
due analitiche (di cui una cross-database). La Tabella~\ref{tab:query-class}
mostra come le nostre sei query coprono e superano il requisito.

\begin{table}[h]
\centering
\small
\begin{tabular}{llll}
\toprule
\textbf{ID} & \textbf{Tipo} & \textbf{Database} & \textbf{Descrizione} \\
\midrule
Q1 & lookup    & Neo4j     & Transfer uscenti da un conto in una finestra temporale \\
Q2 & lookup    & MongoDB   & Profilo completo di una persona (anagrafica + conti + prestiti) \\
Q3 & lookup    & Cross-DB  & Scheda conto + intestatario (Mongo), top-10 vicini (Neo4j) \\
Q4 & analitica & Neo4j     & Cicli di transfer $A\to B\to C\to A$ sopra soglia \\
Q5 & analitica & MongoDB   & Top-10 conti per volume trasferito in un mese \\
Q6 & analitica & Cross-DB  & Hub per grado, soglia o top-$x\%$ (Neo4j) + intestatario (Mongo) \\
\bottomrule
\end{tabular}
\caption{Le sei query: tre di lookup e tre analitiche, due delle quali cross-database.}
\label{tab:query-class}
\end{table}

\paragraph{Esempio di query a grafo (Q4, cicli).} La forza del modello a grafo
si vede nella query dei cicli: un pattern che in SQL sarebbe un incubo di
self-join, in Cypher \`e una sola riga di \emph{pattern matching}.

\begin{lstlisting}[style=plain,caption={Q4 -- ricerca di cicli a tre passi in Cypher.},captionpos=b]
MATCH (a:Account)-[t1:TRANSFER]->(b:Account)
      -[t2:TRANSFER]->(c:Account)-[t3:TRANSFER]->(a)
WHERE t1.amount >= $minAmount AND t2.amount >= $minAmount
  AND t3.amount >= $minAmount
  AND t1.timestamp < t2.timestamp
  AND t2.timestamp < t3.timestamp
  AND t3.timestamp - t1.timestamp <= $windowMs
RETURN a.id, b.id, c.id, t1.amount, t2.amount, t3.amount
LIMIT 50
\end{lstlisting}

\paragraph{Esempio di query cross-database (Q6).} Q6 fa lavorare i due motori in
sequenza: prima Neo4j trova gli hub (i conti con grado alto, formula
\ref{eq:degree}), poi MongoDB arricchisce ciascun hub con tipo di conto, data
di apertura e profilo dell'intestatario (persona o azienda). \`E il pattern
``il grafo trova \emph{chi}, il documento dice \emph{cosa}''.

\begin{lstlisting}[style=py,caption={Q6 -- la parte di orchestrazione cross-database (Python).},captionpos=b]
# 1) Neo4j: trova gli hub per grado dentro la finestra
#    (CYPHER_THRESHOLD se si usa la soglia, CYPHER_TOP_PCT se si usa il top x%)
hubs = session.run(CYPHER_THRESHOLD, startMs=..., endMs=..., minDegree=...).data()
ids  = [h["accountId"] for h in hubs]

# 2) MongoDB: metadati del conto + intestatario, con una find($in) per collezione
meta   = {d["_id"]: d for d in db.account.find({"_id": {"$in": ids}})}
owners = owners_of(ids)   # legge personOwnAccount e companyOwnAccount
\end{lstlisting}

L'helper \texttt{owners\_of} (in \texttt{app/queries/\_common.py}) \`e condiviso
con Q3: risolve l'intestatario passando dalle due collezioni-link
(\texttt{personOwnAccount}, \texttt{companyOwnAccount}) alle rispettive
anagrafiche, sempre con query \texttt{\$in} in blocco e mai una query per conto.

\paragraph{Esempio di aggregazione (Q5, MongoDB).} Dove Neo4j usa il pattern
matching, MongoDB usa la \emph{aggregation pipeline}: una sequenza di stadi che
trasformano i documenti a catena. Q5 trova i dieci conti con pi\`u volume
trasferito in un mese.

\begin{lstlisting}[style=py,caption={Q5 -- aggregation pipeline di MongoDB.},captionpos=b]
pipeline = [
    {"$match":  {"timestamp": {"$gte": start_ms, "$lt": end_ms}}},
    {"$group":  {"_id": "$accountSrc",
                 "totalAmount":  {"$sum": "$amount"},
                 "numTransfers": {"$sum": 1}}},
    {"$sort":   {"totalAmount": -1}},
    {"$limit":  10},
    {"$lookup": {"from": "account", "localField": "_id",
                 "foreignField": "_id", "as": "account"}},
]
\end{lstlisting}

Si legge come una catena di trasformazioni: \texttt{\$match} filtra il mese
(lo stadio selettivo che rende la query veloce), \texttt{\$group} somma gli
importi per conto, \texttt{\$sort} e \texttt{\$limit} tengono i primi dieci,
\texttt{\$lookup} aggancia i metadati del conto. \`E l'equivalente
document-oriented di un \texttt{GROUP BY} con join. Avere i due paradigmi a
confronto --- pattern matching dichiarativo su grafo e pipeline procedurale su
documenti --- \`e uno dei punti pi\`u istruttivi del progetto.

\subsection{Interfaccia web}

L'interfaccia (Flask + Tailwind) presenta le sei query come schede selezionabili;
per ogni query mostra un form con i parametri. Una funzione di comodo
(\texttt{/api/random}) pesca valori \emph{realmente presenti} nel dataset, cos\`i
chi prova l'app non deve conoscere a memoria gli identificativi dei conti (che
sono interi a 19 cifre). Questa scelta nasce da un problema concreto di
usabilit\`a: senza, l'utente non saprebbe cosa scrivere nei campi e otterrebbe
sempre zero risultati.

% ===================================================================
\section{Risultati e analisi critica}
\label{sec:risultati}
% ===================================================================

Abbiamo misurato ogni query eseguendola cinque volte di fila con parametri di
default realistici. La Tabella~\ref{tab:perf} riporta il numero di righe
restituite e il tempo \emph{a regime} (la mediana delle esecuzioni), insieme al
tempo della \emph{prima} esecuzione, sempre pi\`u alto.

\begin{table}[h]
\centering
\begin{tabular}{lllrrr}
\toprule
\textbf{ID} & \textbf{Tipo} & \textbf{DB} & \textbf{Righe} & \textbf{A regime (ms)} & \textbf{Prima run (ms)} \\
\midrule
Q1 & lookup    & Neo4j    & 100 & 10,0   & 169,7 \\
Q2 & lookup    & MongoDB  & 1   & 0,5    & 63,9 \\
Q3 & lookup    & Cross-DB & 1$^{*}$ & 6,2 & 67,4 \\
Q4 & analitica & Neo4j    & 50  & 660,8  & 906,9 \\
Q5 & analitica & MongoDB  & 10  & 1,0    & 35,9 \\
Q6 & analitica & Cross-DB & 20  & 1289,6 & 1915,1 \\
\bottomrule
\end{tabular}
\caption{Prestazioni misurate (5 esecuzioni). $^{*}$Q3 restituisce un documento
che contiene a sua volta i 10 conti vicini.}
\label{tab:perf}
\end{table}

I numeri raccontano una storia molto netta e coerente con la teoria della
Sezione~2.

\paragraph{Q1 --- Transfer uscenti (lookup, Neo4j).} 100 righe in circa 10~ms.
La query \`e molto selettiva: parte da un singolo conto (indice su
\texttt{Account.id}) e segue solo i suoi archi uscenti dentro la finestra. \`E il
caso ideale per il grafo: navigazione locale a partire da un nodo noto.

\paragraph{Q2 --- Profilo persona (lookup, MongoDB).} Un solo documento in
mezzo millisecondo, la query pi\`u veloce in assoluto. La persona \`e trovata per
\texttt{\_id} (indice automatico) e il documento contiene gi\`a, tramite quattro
\emph{lookup} a catena, i suoi conti e prestiti. \`E la dimostrazione del
vantaggio del modello a documenti per le schede aggregate.

\paragraph{Q3 --- Scheda conto + vicini (lookup, cross-DB).} Circa 6~ms. Qui i
due database collaborano: MongoDB d\`a i metadati del conto e dell'intestatario,
Neo4j i dieci conti pi\`u connessi per numero di transfer. Il costo \`e ancora
basso perch\'e tutto parte da un singolo conto noto.

\paragraph{Q4 --- Cicli a tre passi (analitica, Neo4j).} Qui il tempo sale a
circa 660~ms, due ordini di grandezza sopra le lookup. \`E atteso: la ricerca
del pattern $A\to B\to C\to A$ \`e poco selettiva, perch\'e il motore deve
esplorare moltissimi cammini di lunghezza tre prima di scoprire quali si chiudono
in ciclo. \`E il prezzo di un'interrogazione potente che sul relazionale sarebbe
ancora pi\`u cara.

\paragraph{Q5 --- Top-10 volume mensile (analitica, MongoDB).} Circa 1~ms,
sorprendentemente veloce per un'analitica. Il motivo \`e che l'aggregazione opera
su un filtro temporale selettivo (un solo mese) e MongoDB la risolve con la sua
pipeline nativa. Mostra che ``analitica'' non significa automaticamente ``lenta'':
conta quanto \`e selettivo il filtro a monte.

\paragraph{Q6 --- Hub detection (analitica, cross-DB).} La query pi\`u pesante,
circa 1,3~secondi. Deve scorrere \emph{tutti} i transfer dentro la finestra
temporale, calcolare il grado di ogni conto (formula~\ref{eq:degree}) e tenere
solo quelli sopra soglia. \`E poco selettiva per costruzione --- guarda l'intera
rete --- e per questo \`e la pi\`u costosa. La variante top-$x\%$ costa ancora
un po' di pi\`u (circa 1,5~s a regime con $x=1$): non basta pi\`u tenere i
primi 20 in un \emph{heap}, bisogna ordinare e materializzare tutta la
classifica per sapere dove cade il taglio percentuale. Un dettaglio
implementativo che ha contato: raggruppare per \emph{nodo} e leggere
\texttt{a.id} solo sui conti sopravvissuti al filtro, invece che su tutti i
2,7 milioni di estremi di arco, vale circa 0,7~s.

\paragraph{Il costo della ``prima volta''.} In tutte le query la prima
esecuzione \`e nettamente pi\`u lenta delle successive (per Q2 si passa da 64~ms
a mezzo millisecondo). \`E l'effetto del \emph{caching}: alla prima esecuzione i
dati non sono ancora nella memoria centrale dei database e vanno letti da disco;
dopo, restano nel buffer e le query a regime sono molto pi\`u veloci. \`E la
\emph{localit\`a temporale} descritta negli appunti (cap. 2, gestione del buffer):
un dato appena letto ha alta probabilit\`a di essere riletto a breve, quindi
tenerlo in memoria ammortizza il costo.

\paragraph{Lettura d'insieme.} Le quattro lookup/analitiche selettive (Q1, Q2,
Q3, Q5) stanno tutte sotto i 10~ms; le due analitiche di grafo poco selettive
(Q4, Q6) costano centinaia/migliaia di ms. La discriminante non \`e ``grafo
contro documento'' n\'e ``lookup contro analitica'', ma la \textbf{selettivit\`a}
(formula~\ref{eq:sel}): pi\`u dati la query \`e costretta a toccare, pi\`u tempo
impiega. \`E esattamente la lezione del capitolo sull'ottimizzazione.

% ===================================================================
\section{Limitazioni}
% ===================================================================

\begin{itemize}[nosep]
  \item \textbf{Ridondanza e coerenza.} Tenere i transfer in due database
  significa che un aggiornamento andrebbe propagato a entrambi. Nel prototipo i
  dati sono in sola lettura, quindi il problema non si pone; in un sistema reale
  servirebbe un meccanismo di sincronizzazione (es. \emph{change data capture}).
  \item \textbf{Scala singola.} Tutto gira su un nodo per database. Q6, che
  scorre l'intera rete, a scale factor molto pi\`u alti diventerebbe proibitiva e
  richiederebbe sharding o pre-calcolo dei gradi (materializzazione).
  \item \textbf{Hub detection: due criteri, nessuno ``giusto''.} Q6 offre
  soglia assoluta e top-$x\%$, ma la scelta della soglia (o della percentuale)
  resta a carico dell'analista: un criterio pi\`u robusto userebbe statistiche
  della distribuzione (es. media pi\`u $k$ deviazioni standard, o percentili
  calcolati per periodo).
  \item \textbf{Nessuna elaborazione a flusso.} Il sistema lavora su dati
  statici. Un caso antifrode reale, che deve reagire alle transazioni in tempo
  reale, richiederebbe stream processing (Sezione~\ref{sec:distr}).
\end{itemize}

% ===================================================================
\section{Conclusioni}
% ===================================================================

FinDetect mostra in pratica perch\'e la polyglot persistence \`e una scelta
sensata su dati finanziari: le interrogazioni sui cammini (cicli, hub) vivono
naturalmente nel grafo, le schede aggregate nel documento, e ogni motore lavora
dove rende di pi\`u. L'analisi dei tempi conferma la teoria del corso: il costo
di una query \`e governato dalla sua selettivit\`a, non dal tipo di database. Le
sei query coprono e superano il requisito della traccia (lookup e analitiche,
con casi cross-database), e l'interfaccia rende il tutto dimostrabile senza
conoscere a memoria gli identificativi del dataset. I limiti principali ---
ridondanza, scala singola, assenza di stream --- sono noti e indicano le
naturali direzioni di sviluppo.

% ===================================================================
\section{Comandi per l'esecuzione}
% ===================================================================

\begin{lstlisting}[style=plain,caption={Avvio completo del prototipo.},captionpos=b]
# 1) database in Docker (Neo4j + MongoDB replica set)
docker compose up -d

# 2) generazione del dataset SF 1 (una volta sola, ~8 min)
./datagen/run_datagen.sh

# 3) ambiente Python
python3.11 -m venv .venv
source .venv/bin/activate
pip install -r requirements.txt

# 4) caricamento dati nei due database
python -m etl.load_mongo
python -m etl.load_neo4j

# 5) avvio dell'app web -> http://localhost:5050
python -m app.app

# 6) (opzionale) smoke test: le 6 query coi parametri di default
python test_smoke.py
\end{lstlisting}

\end{document}
